\chapter*{Словарь терминов}             % Заголовок
\addcontentsline{toc}{chapter}{Словарь терминов}  % Добавляем его в оглавление

\textbf{Автономный мобильный робот} --- мобильная робототехническая платформа, способная перемещаться в рабочей среде, воспринимать окружение с помощью датчиков, строить или использовать карту, планировать маршрут и выполнять задачи без жёсткой привязки к физической направляющей инфраструктуре.

\textbf{Роботизированный склад} --- складская система, в которой значительная часть операций перемещения, комплектации, сортировки или транспортировки поддерживается автономными мобильными роботами и программной системой управления.

\textbf{Периферийные вычисления} --- архитектура размещения вычислительных ресурсов вблизи конечных устройств, датчиков, роботов и локальных точек беспроводного доступа для снижения задержек и разгрузки удалённого облака.

\textbf{Периферийный узел} --- локальный вычислительный ресурс, размещённый вблизи источников данных и предназначенный для обработки задач с меньшей задержкой по сравнению с удалённым облаком.

\textbf{Вычислительная разгрузка} --- передача части вычислений с устройства-источника задачи, например автономного мобильного робота, на внешний вычислительный узел, расположенный на периферии сети или в облаке.

\textbf{Управление допуском задач} --- процедура принятия решения о том, какие поступающие задачи следует принять к обслуживанию, отложить, перенаправить или отклонить при ограниченных вычислительных ресурсах.

\textbf{Проектирование экономических механизмов} --- раздел теории игр и экономической теории, изучающий построение правил взаимодействия участников, при которых их локальные решения приводят к требуемому системному результату.

\textbf{Теневая цена} --- внутренняя оценка дефицитности ресурса или внешнего эффекта использования ресурса, не обязательно являющаяся реальным денежным платежом.

\textbf{Тип участника механизма} --- локальная информация участника, существенная для распределения ресурсов: для вычислительной задачи это класс, срок, приоритет, ресурсные требования и ценность; для периферийного узла --- доступная ёмкость, очередь, состояние отказа и параметры стоимости обслуживания.

\textbf{Структурированная заявка} --- сообщение участника механизма, содержащее параметры, необходимые для вычисления полезности, стоимости и допустимости назначения. В системе одного оператора такая заявка не является денежной ставкой в узком смысле.

\textbf{Экспонируемая ёмкость} --- часть физически свободных ресурсов узла, предоставляемая для распределения новых задач в текущем раунде с учётом выбранного операционного режима.

\textbf{Внешний эффект назначения задачи} --- изменение достижимого качества обслуживания остальных задач, возникающее из-за того, что рассматриваемая задача использует часть общего дефицитного ресурса.

\textbf{Общественное благосостояние} --- критерий качества распределения, равный сумме чистых вкладов назначенных задач, где чистый вклад определяется как разность между полезностью выполнения задачи и стоимостью использования периферийного узла.

\textbf{Задача выбора победителей} --- комбинаторная задача определения допустимого распределения задач между вычислительными узлами, максимизирующего выбранный критерий качества, например общественное благосостояние.

\textbf{Операционный режим периферийного узла} --- дискретное состояние локального управления узлом, определяющее долю ресурсов, доступную для новых задач, правила резервирования ёмкости и приоритеты обслуживания.

\textbf{VCG-платёж} --- внутренняя цена, вычисляемая по правилу Викри--Кларка--Гровса и отражающая внешний эффект участника на остальных участников механизма.

\textbf{Нулевая альтернатива} --- допустимый исход, соответствующий неучастию задачи или участника в распределении, например отказу от назначения или локальному выполнению задачи.

\textbf{Локальная эффективность аукционного раунда} --- свойство механизма выбирать в данном раунде распределение, максимизирующее общественное благосостояние среди допустимых вариантов при фиксированном состоянии системы.

\textbf{Стратегическая совместимость} --- свойство механизма, при котором участнику невыгодно искажать свою оценку при прочих равных условиях.

\textbf{Индивидуальная рациональность} свойство механизма, при котором участие в нём не ухудшает положение участника по сравнению с допустимой альтернативой неучастия.

\textbf{Аукционный слой распределения} --- компонент метода, который на основе заявок участников, доступных ресурсов и критерия общественного благосостояния определяет назначение задач и соответствующие платежи.

\textbf{Централизованное обучение и децентрализованное исполнение} --- подход к многоагентному обучению, при котором на этапе обучения используется глобальная информация о системе, а на этапе исполнения каждый агент принимает решение по локальному наблюдению.

\textbf{Командное вознаграждение} --- общий сигнал качества, одинаковый для всех обучаемых агентов и отражающий системную цель, а не локальную выгоду отдельного узла.

\textbf{Инварианты разделения контуров} --- ограничения, согласно которым обучаемая политика не изменяет заявки роботов, функции полезности, функции стоимости и правило аукционного распределения, а управляет только операционными режимами периферийных узлов.

\textbf{Аукционные признаки} --- показатели, формируемые аукционным слоем после распределения задач: платежи, оценки внешнего эффекта, вклад в общественное благосостояние, доля использованной экспонируемой ёмкости и доля отклонённых задач.

\textbf{Кластерный аукцион} --- локальный аукционный раунд, выполняемый внутри подграфа периферийной инфраструктуры для уменьшения вычислительных и коммуникационных накладных расходов.

\textbf{Межкластерный резерв} --- ограниченный набор периферийных узлов соседних кластеров, которые могут использоваться для задач перегруженного кластера с дополнительной сетевой стоимостью.

\textbf{Метод AMPLE-AMR } --- предлагаемый в диссертации гибридный метод адаптивного распределения вычислительных задач автономных мобильных роботов. Метод сочетает обучаемую многоагентную политику управления режимами периферийных узлов и VCG-подобный аукционный механизм распределения задач в рамках текущих ресурсных ограничений.

\textbf{Метод C-AMPLE-AMR} --- масштабируемое расширение метода AMPLE-AMR, использующее кластеризацию графа периферийной инфраструктуры и проведение локальных аукционов внутри выделенных кластеров.

\textbf{Доменная параметризация} --- задание параметров прикладной среды, необходимых для воспроизводимой проверки метода в выбранном предметном домене.

\textbf{Профиль нагрузки} --- правило изменения интенсивности и состава вычислительных задач во времени, используемое для моделирования стабильного, пикового, смешанного или аварийного режима работы системы.

\textbf{Зона ухудшенной связи} --- область складской среды, в которой увеличиваются задержки передачи данных, снижается пропускная способность или возрастает вероятность временной потери соединения.

\textbf{Профиль масштабирования} --- экспериментальный профиль, в котором число роботов, задач и периферийных узлов увеличивается для оценки вычислительных и коммуникационных накладных расходов.
